\title{FlashAttention-3: Fast and Accurate Attention with Asynchrony and Low-precision}

\begin{abstract}
Attention mechanisms, integral to the Transformer architecture that underpins most modern language models, represent the principal computational bottleneck, particularly in scaling to larger models or handling extended contexts. \textsc{FlashAttention} previously demonstrated how optimizing memory access patterns could accelerate attention computation on GPUs, yet this method does not fully exploit the enhanced functionalities found in newer hardware. Indeed, \textsc{FlashAttention-2} utilizes only 35\% of the available throughput on the H100 GPU. To address this limitation, we introduce three primary strategies tailored for the Hopper GPU platform: leveraging asynchronous capabilities of Tensor Cores and TMA to (1) concurrently mask computation latency and data transfer via warp-specialized execution, (2) merge block-wise matrix multiplication with the softmax stage, and (3) employ block-level quantization and out-of-order execution via hardware-accelerated FP8 precision. Our proposed technique, \textsc{FlashAttention-3}, delivers a 1.5–2.0$\times$ improvement in speed on H100 GPUs, with FP16 performance peaking at 740 TFLOPs/s (equating to 75\% utilization), and FP8 approaching nearly 1.2 PFLOPs/s. In addition, we confirm that FP8 \textsc{FlashAttention-3} exhibits 2.6$\times$ less numerical inaccuracy than conventional FP8 attention methods.
\end{abstract}

\section{Introduction}
\label{sec:intro}

Within the Transformer architecture~\citep{vaswani2017attention}, the attention mechanism represents a central computational constraint, as calculating self-attention scores between queries and keys requires operations that scale quadratically with the sequence length. Enabling attention to process longer contexts opens up a variety of novel capabilities, such as integrating and reasoning across multiple extensive documents~\citep{guo2021longt5,shaham2022scrolls,peng2023yarn}, analyzing files within extensive codebases~\citep{roziere2023code, li2023starcoder}, and accommodating new data modalities like high-resolution imagery~\citep{chen2022scaling}, audio~\citep{gulati2020conformer}, and video~\citep{ho2022video}. Furthermore, it supports emerging use cases, including managing user interactions with prolonged history~\citep{sun2019bert4rec} and guiding agent workflows over extended time horizons~\citep{yao2022react}. This challenge has motivated efforts to accelerate attention mechanisms in settings with long context sequences, employing approaches such as approximate methods~\citep{katharopoulos2020transformers,choromanski2020rethinking, tay2020efficient}, advances in software optimization~(\citep{rabe2021self,dao2022flashattention,kwon2023efficient}), and even designing alternative model architectures~\citep{peng2023rwkv,sun2023retentive,gu2023mamba}.

This study builds upon the methodology proposed by \citet{dao2022flashattention}, which focuses on designing exact-attention algorithms that account for the GPU execution model and its architectural nuances within their overall approach. Dao et al., in \citep{dao2022flashattention}, introduced \textsc{FlashAttention}—a groundbreaking tiling method that leverages parallelism by integrating all attention-related computations into a single GPU kernel, thereby avoiding unnecessary intermediate accesses to slow global memory. The subsequent work by \citet{dao2023flashattention2} modified the algorithm into \textsc{FlashAttention-2}, facilitating parallelization along the sequence length axis and organizing the innermost loop of the forward computation to operate over partitions of the key and value matrices. This adjustment further enhances workload allocation and GPU utilization. Despite these optimizations, we note that \textsc{FlashAttention-2} still demonstrates suboptimal performance on the latest GPU hardware compared to highly-tuned matrix multiplication (GEMM) kernels (achieving only 35\% utilization versus 80–90\% on the Hopper H100 GPU). Some of these discrepancies are likely linked to implementation-level choices, such as utilizing Ampere instructions rather than Hopper-optimized commands on Tensor Cores when running on Hopper architecture. Recent advancements, including ThunkerKitten~\citep{spector2024thunder} and cuDNN 9~\citep{cudnn9}, provide evidence that the adoption of Hopper-native instructions and tile-centric abstractions can both accelerate attention computation and streamline code development.

At a more foundational level, the algorithm underlying \textsc{FlashAttention-2} is constructed around a basic synchronous execution model, deliberately omitting direct utilization of asynchronous processing and low-precision arithmetic within its framework. Asynchrony stems from advances in hardware specialization, where critical operations in machine learning workloads are accelerated through dedicated units—such as Tensor Cores for matrix multiplications and the Tensor Memory Accelerator (TMA) for memory access—which operate independently from the broader array of CUDA cores handling logical, integer, and floating-point tasks. Meanwhile, support for reduced-precision computation, exemplified by FP8 on Hopper and FP4 on Blackwell (following the precedent set by FP16 in Pascal [2017] and BF16 in Ampere [2020]), has consistently delivered twofold or fourfold improvements in throughput for equivalent energy consumption and silicon real estate. The pertinent features of the Hopper architecture concerning these aspects are discussed in detail in \cref{subsec:hardware}. The principal engineering obstacle lies in revising \textsc{FlashAttention-2} so that it can effectively exploit these hardware mechanisms: leveraging asynchrony means devising ways to overlap tasks such as matmul and softmax, even where dependency chains exist, while integrating low-precision requires careful strategies to control quantization artifacts—an issue heightened when handling outlier activations in large language models~\citep{dettmers2208llm, sun2024massive}.

In pursuit of this goal, we introduce \textsc{FlashAttention-3}, which integrates and advances three novel concepts aimed at enhancing efficiency on modern GPU hardware.\footnote{While our findings focus primarily on NVIDIA's Hopper architecture, our method generalizes to any GPU platform offering strong asynchronous processing and support for low-precision operations.}

\begin{enumerate}

\item \textbf{Producer-Consumer asynchrony:} We propose a warp-specific software pipelining strategy that leverages the ability of data movement and Tensor Core computations to proceed asynchronously. By segregating the production and consumption roles into different warps, the method enhances latency hiding with respect to both memory access and instruction issuing.
\item \textbf{Concealing softmax within asynchronous block-wise GEMMs:} Non-GEMM softmax-related computations, such as floating point multiply-add and exponentials, are scheduled concurrently with asynchronous WGMMA GEMM instructions. To facilitate this, we modify the \textsc{FlashAttention-2} algorithm so as to avoid certain serial dependencies linking softmax operations and GEMMs. For instance, in our algorithm’s 2-stage variant, softmax is processed on one block of the scores matrix, while WGMMA initiates asynchronously to prepare the next block.
\item \textbf{Low-precision GEMM leveraging hardware acceleration:} Our forward algorithm is adapted to utilize FP8 Tensor Cores for GEMM, achieving nearly a twofold increase in observed TFLOPs/s. This adaptation addresses WGMMA’s particular requirements concerning memory layouts of FP32 accumulators and FP8 operand matrices. To counteract the precision loss incurred by operating in FP8, we employ block quantization and incoherent processing techniques.
\end{enumerate}

For empirical assessment, we evaluated the performance of \textsc{FlashAttention-3} on an H100 SXM5 GPU, spanning various parameter configurations. The results demonstrate that (1) using FP16 yields a speed increase of 1.5-2.0$\times$ over \textsc{FlashAttention-2} during the forward computation (reaching up to 740 TFLOPs/s), and offers a 1.5-1.75$\times$ boost in the backward computation; (2) the FP8 variant nearly attains 1.2 PFLOPs/s; and (3) at extensive sequence lengths, FP16 consistently performs better, while FP8 shows comparable performance\footnote{Specifically, FP8 surpasses for head dimension 64, and matches performance for head dimensions 128 and 256 without causal masking, but lags in cases with causal masking.} relative to the advanced attention routines found in NVIDIA’s cuDNN library. Additionally, the FP16 version of \textsc{FlashAttention-3} maintains the numerical error observed with \textsc{FlashAttention-2}, and improves upon standard attention implementations by preserving intermediate calculations (such as softmax rescaling) in FP32. Notably, FP8 \textsc{FlashAttention-3}, utilizing block quantization along with incoherent processing, achieves 2.6$\times$ higher accuracy compared to standard attention employing per-tensor quantization, particularly in scenarios with outlier features.

We have released \textsc{FlashAttention-3} as open-source under a permissive license\footnote{\textsc{FlashAttention-3} is available at \texttt{https://github.com/Dao-AILab/flash-attention}} and intend to incorporate it into the PyTorch and Hugging Face ecosystems, thereby maximizing its accessibility to the research and development community.

\section{Background: Multi-Head Attention and GPU Characteristics}
\label{sec:background}

\subsection{Multi-Head Attention}
\label{subsec:multi_head_attn}

Let $\mathbf{Q}, \mathbf{K}, \mathbf{V} \in \mathbb{R}^{N \times d}$ be the query, key and value input sequences associated to a single head, where $N$ is the sequence length and $d$ is the head dimension. Then the attention output $\mathbf{O}$ is computed as:

\begin{equation*}
  \mathbf{S} = \alpha \mathbf{Q} \mathbf{K}^\top \in \mathbb{R}^{N \times N}, \quad \mathbf{P} = \mathrm{softmax}(\mathbf{S}) \in \mathbb{R}^{N \times N}, \quad \mathbf{O} = \mathbf{P}\mathbf{V} \in \mathbb{R}^{N \times d},
\end{equation*}

The $\mathrm{softmax}$ function is computed along each row, and it is common to choose $\alpha = 1/\sqrt{d}$ to scale the inputs. To mitigate numerical issues arising from the exponentiation, $\mathrm{rowmax}(\mathbf{S})$ is subtracted from $\mathbf{S}$ before applying the softmax. In the context of multi-head attention (MHA), individual heads operate with separate query, key, and value projections, allowing for simultaneous computation across several heads and batches to generate the complete output tensor.

Now let $\phi$ be a scalar loss function and let $\mathbf{d}(-) = \partial \phi / \partial (-)$ be notation for the gradient. Given the output gradient $\mathbf{dO} \in \mathbb{R}^{N \times d}$, we compute $\mathbf{dQ}$, $\mathbf{dK}$, and $\mathbf{dV}$ according to the chain rule as follows:

\begin{align*}
  \mathbf{dV} &= \mathbf{P}^\top \mathbf{dO} \in \mathbb{R}^{N \times d} \\
  \mathbf{dP} &= \mathbf{dO} \mathbf{V}^\top \in \mathbb{R}^{N \times N} \\
  \mathbf{dS} &= \mathrm{dsoftmax} (\mathbf{dP}) \in \mathbb{R}^{N \times N} \\
  \mathbf{dQ} &= \alpha \mathbf{dS} \mathbf{K} \in \mathbb{R}^{N \times d} \\
  \mathbf{dK} &= \alpha \mathbf{dS}^\top \mathbf{Q} \in \mathbb{R}^{N \times d},
\end{align*}

In this context, we express $\mathbf{d}s = (\mathrm{diag}(p) - p p^\top)\mathbf{d}p$ where $p = \mathrm{softmax}(s)$ with respect to a vector $s$, using $\mathrm{dsoftmax}(\mathbf{dP})$ to denote applying this operation to each row individually. Furthermore, during the backward pass of MHA, this calculation can be efficiently executed in parallel over the set of heads and batches.

\subsection{GPU hardware characteristics and execution model}
\label{subsec:hardware}

We outline key features of the GPU execution model pertinent to \textsc{FlashAttention-3}, specifically emphasizing the NVIDIA Hopper architecture to provide a practical example of this framework.

\paragraph{Memory hierarchy:} On GPUs, memory is structured as a layered set of storage regions, where memory bandwidth increases as capacity decreases (\cref{tab:gpu-hierarchy})\footnote{\citet{luo2024benchmarking} indicates that shared memory offers a bandwidth of 128 bytes per clock cycle per SM; to estimate total bandwidth, we multiply by 132 SMs and the boost clock rate of 1830 MHz.}. Global memory (GMEM), often referred to as HBM, is an external DRAM available to all streaming multiprocessors (SMs). Information residing in GMEM is automatically placed into the on-chip L2 cache. Additionally, every SM has access to a compact, highly banked, programmer-controlled cache known as shared memory (SMEM). The innermost component in this hierarchy is the register file unique to each SM.

\paragraph{Thread hierarchy:} In the GPU programming paradigm, execution is structured into logical groups known as threads. These are arranged in a hierarchy, ranging from the smallest unit up to the largest: individual threads, followed by warps (which consist of 32 threads), then warpgroups (aggregating 4 sequential warps), progressing to threadblocks (also referred to as cooperative thread arrays or CTAs), threadblock clusters (a construct introduced with Hopper), and, at the highest level, grids.

The relationship between these two hierarchies is highly interdependent. Threads belonging to a single CTA are simultaneously scheduled on the identical SM, while CTAs within the same cluster are executed together on the same GPC. All threads of a CTA have direct access to the SMEM, but each thread is limited to a maximum of 256 private registers (RMEM) that are exclusively accessible by that thread.

\paragraph{Asynchrony and warp-specialization:}

GPUs function as throughput-oriented processors that utilize parallelism and asynchronous operations to mask delays associated with memory accesses and computation. To facilitate asynchronous memory transfers between GMEM and SMEM, Hopper introduces the Tensor Memory Accelerator (TMA), a specialized hardware device \cite[\S7.29]{cuda}. Additionally, Hopper departs from previous designs like Ampere by enabling its Tensor Core—accessible through the warpgroup-wide WGMMA instruction \cite[\S9.7.14]{ptx}—to operate asynchronously and to fetch its input data directly from shared memory.

Hardware mechanisms for supporting asynchrony enable the design of warp-specialized kernels, where warps within a CTA are assigned distinct functions, acting exclusively as either producers (handling data transfers) or consumers (performing computations). This architectural feature generally enhances the compiler's capacity to produce efficient instruction schedules \citep{warp-specialization-2011}. Furthermore, Hopper provides flexibility in register allocation among warpgroups by means of \verb|setmaxnreg| \cite[\S9.7.17.1]{ptx}, which permits warps engaged in MMAs to access a greater portion of RMEM, in contrast to those responsible only for issuing TMA—operations typically requiring only a single thread.

\paragraph{Low-precision number formats:}
\label{sec:low-precision-gpu}
Contemporary GPUs are equipped with dedicated hardware designed to enhance the speed of low-precision arithmetic. For instance, the WGMMA instruction leverages FP8 Tensor Cores on Hopper, enabling each SM to achieve twice the computational throughput relative to operations performed in FP16 or BF16.

Nevertheless, employing FP8 WGMMA correctly requires a clear understanding of the operand layout restrictions. Consider a GEMM operation that computes $A \times B^{\top}$, where $A$ is of size $M \times K$ and $B$ is of size $N \times K$. We define the $A$ or $B$ operand as \emph{mn-major} if the outermost $M$ or $N$ dimension is stored contiguously, or as \emph{k-major} if the innermost $K$ dimension is contiguous in memory. While FP16 WGMMA is flexible, accepting both mn-major and k-major formats for shared memory operands, FP8 WGMMA imposes a stricter requirement: only k-major formats are permitted. This becomes particularly challenging in contexts like attention mechanisms, where chaining multiple GEMMs within a single kernel is desirable. In such cases, incompatibilities between the layouts of FP32 accumulators and FP8 operands hinder the invocation of successive FP8 WGMMAs.

With respect to attention mechanisms, these layout constraints require specific alterations to how the FP8 algorithm is structured, as outlined in detail in \cref{sec:algofp8}.

\subsection{Standard Attention and Flash Attention} As described by \citet{dao2022flashattention}, \textbf{standard attention} refers to the GPU-based implementation of attention in which the intermediate matrices $\mathbf{S}$ and $\mathbf{P}$ are explicitly stored in HBM. \textsc{FlashAttention} introduces a key optimization: it applies a localized form of the softmax reduction, thereby minimizing costly memory accesses required to write and read intermediate data, and consolidating attention computation within a unified kernel. The notion of local softmax maps to lines \ref{code-ws:softmax_start}-\ref{code-ws:softmax_end} in the consumer mainloop of \cref{alg:flash3_wgmma_ws_only}, in combination with the normalization of $\mathbf{O}$ block segments. For a clear and succinct argument that this strategy correctly produces $\mathbf{O}$, see \cite[\S 2.3.1]{dao2023flashattention2}.

\section{FlashAttention-3: Algorithm}
\label{sec:algo}

This section presents an overview of the \textsc{FlashAttention-3} algorithm. Our exposition centers on the forward pass; the procedure for the backward pass is outlined in~\cref{sec:algo_ws_bwd}. Initially, we demonstrate the incorporation of warp-specialization and a circular SMEM buffer into the foundational framework of \textsc{FlashAttention-2}. Next, we elaborate on leveraging WGMMA’s asynchronous characteristics to construct a two-stage pipeline where GEMM and softmax operations are overlapped. We conclude by detailing the adaptations required for FP8—addressing both layout compliance and accuracy concerns through block quantization and incoherent computation.

\subsection{Producer-Consumer asynchrony through warp-specialization and pingpong scheduling}
\label{sec:algo_ws}

\paragraph{Warp-specialization}
Much like in \textsc{FlashAttention-2}, the forward computation of \textsc{FlashAttention-3} can be highly parallelized across batch dimensions, attention heads, and the length of the query sequence. Therefore, it is adequate to present the algorithm at the CTA level, where a tile $\mathbf{Q}_i$ from the query matrix is used to compute its corresponding output tile $\mathbf{O}_i$. To make the explanation clearer, we initially introduce the warp-specialized approach employing a circular SMEM buffer, intentionally excluding any overlapping of GEMM and softmax operations. Suppose $d$ represents the head dimension and $N$ is the total sequence length; then, by setting the query block size to $B_r$, we partition $\mathbf{Q}$ into $T_r = \lceil \frac{N}{B_r} \rceil$ distinct blocks, namely $\mathbf{Q}_1, .., \mathbf{Q}_{T_r}$.

\begin{algorithm}[H]
    \caption{\small\label{alg:flash3_wgmma_ws_only}\textsc{FlashAttention-3} Forward Computation \textbf{Excluding} Intra-Consumer Overlap -- CTA Perspective}
    \begin{algorithmic}[1]
\REQUIRE Input matrices $\mathbf{Q}_i \in \mathbb{R}^{B_r \times d}$, $\mathbf{K}, \mathbf{V} \in \mathbb{R}^{N \times d}$ stored in HBM; key partition size $B_c$ such that $T_c = \lceil \frac{N}{B_c} \rceil$.
\STATE Construct pipeline controller to coordinate barrier synchronization with a $s$-stage cyclic SMEM buffer.
\IF {executing within a producer warpgroup}
\STATE Free a preset number of registers.
\STATE Trigger transfer of $\mathbf{Q}_i$ from HBM into shared memory.
\STATE When done, signal consumer that $\mathbf{Q}_i$ is present in shared memory.
\FOR{$j$ ranging from $0$ to $T_c - 1$}
    \STATE Wait until the $(j\,\%\,s)$ buffer segment is available (consumed).
    \STATE Trigger loads of $\mathbf{K}_j, \mathbf{V}_j$ from HBM into the $(j\,\%\,s)$th portion of shared memory buffer.
    \STATE Once complete, notify consumers that $\mathbf{K}_j, \mathbf{V}_j$ are resident in SMEM.
\ENDFOR
\ELSE 
\STATE Reacquire designated registers according to the current number of consumer warps.
\STATE Allocate on-chip variables: $\mathbf{O}_i = (0) \in \mathbb{R}^{B_r \times d}$ and $\ell_i = 0$, $m_i = -\infty$ in $\mathbb{R}^{B_r}$.
\STATE Wait until $\mathbf{Q}_i$ has been loaded into SMEM.
\FOR{$j$ from $0$ up to $T_c - 1$}
\STATE Await completion of $\mathbf{K}_j$ transfer into shared memory.
\STATE Perform computation $\mathbf{S}_i^{(j)} = \mathbf{Q}_i \mathbf{K}_j^T$ (SS-GEMM). Commit result and pause for synchronization. \label{alg:ws_only_gemm1}
\STATE Save $m_i^{\mathrm{old}} = m_i$ and update $m_i = \mathrm{max}(m_i^{\mathrm{old}}, \mathrm{rowmax}(\mathbf{S}_i^{(j)}))$. \label{code-ws:softmax_start}
\STATE Compute normalized attention weights: $\widetilde{\mathbf{P}}_i^{(j)} = \mathrm{exp}(\mathbf{S}_i^{(j)} - m_i)$, and refresh normalization: $\ell_i = \mathrm{exp}(m_i^{\mathrm{old}} - m_i)\ell_i + \mathrm{rowsum}(\widetilde{\mathbf{P}}_i^{(j)})$. \label{code-ws:softmax_end}
\STATE Wait until $\mathbf{V}_j$ is accessible in SMEM.
\STATE Update output: $\mathbf{O}_i = \mathrm{diag}(\mathrm{exp}(m_i^{\mathrm{old}} - m_i))^{-1} \mathbf{O}_i + \widetilde{\mathbf{P}}_i^{(j)} \mathbf{V}_j$ (RS-GEMM). Commit computation and pause as necessary. \label{alg:ws_only_gemm2}
\STATE Release control of the $(j\,\%\,s)$th buffer segment back to the producer.
\ENDFOR
\STATE Finalize output: $\mathbf{O}_i = \mathrm{diag}(\ell_i)^{-1} \mathbf{O}_i$; store $L_i = m_i + \log(\ell_i)$.
\STATE Transfer $\mathbf{O}_i$ and $L_i$ to HBM as the $i$th entries of $\mathbf{O}$ and $L$.
\ENDIF
\end{algorithmic}
\end{algorithm}

In our implementation of \cref{alg:flash3_wgmma_ws_only} on Hopper, we utilize \verb|setmaxnreg| to manage (de)allocations, employ TMA for loading $\mathbf{Q}_i$ and $\{ \mathbf{K}_j, \mathbf{V}_j \}_{0 \leq j < T_c}$, and apply WGMMA to perform the GEMMs within the consumer mainloop, where the SS or RS prefix indicates the sourcing of the first operand from shared memory or the register file, respectively. To understand the runtime behavior of \cref{alg:flash3_wgmma_ws_only}, it is important to recognize that TMA load operations are asynchronous and therefore not blocked by the completion of previous loads. Furthermore, during the initial $s$ iterations in the producer mainloop, no waiting instructions are issued as the buffer is being populated.

\paragraph{Pingpong scheduling} The asynchronous characteristics of WGMMA and TMA, combined with warp-specialization, present the possibility to simultaneously execute the softmax calculation in one warpgroup while another warpgroup carries out GEMM. To illustrate the motivation, it is important to recognize that matmul operations generally achieve significantly higher throughput than non-matmul computations on current hardware accelerators. For instance, the H100 SXM5 GPU offers a peak performance of 989 TFLOPS for FP16 matmul operations, while its throughput for special functions such as the exponential\footnote{According to the CUDA programming guide, each streaming multiprocessor (SM) is capable of executing 16 special function operations per clock cycle. Using 16, multiplied by 132 SMs and a clock frequency of 1830 MHz, yields 3.9 TFLOPS for special function throughput.} (required for softmax) is only 3.9 TFLOPS. In the attention forward pass with FP16 precision and a head dimension of 128, the number of matmul floating-point operations is 512 times higher than that of exponential operations; however, the exponential function is processed at 256 times lower throughput, leading to the exponential consuming around 50\% of the cycles that matmul operations require. This imbalance becomes even more pronounced in FP8 mode: matmul throughput doubles, yet the exponential throughput remains constant.

Because the exponential operation is handled by a dedicated hardware component known as the multi-function unit, it is preferable to overlap the computation of exponentials with the matmul operations being executed on the Tensor Cores. To achieve this overlap, we utilize synchronization barriers (\texttt{bar.sync} instructions) to control the execution order: we ensure that the GEMMs (GEMM1 — $\mathbf{P} \mathbf{V}$ from the current iteration, and GEMM0 — $\mathbf{Q} \mathbf{K}^\top$ from the subsequent iteration) associated with warpgroup 1 are dispatched prior to those of warpgroup 2. This enables the softmax computation for warpgroup 1 to proceed concurrently with the GEMMs for warpgroup 2. Subsequently, the process alternates so that warpgroup 2 computes the softmax while warpgroup 1 undertakes the GEMMs—a strategy referred to as “pingpong” scheduling, as shown in~\cref{fig:pingpong_scheduling}. Although this pingpong scheduling may not always manifest as precisely as depicted, empirical observations indicate notable performance gains (e.g., throughput rising from 570 TFLOPS to between 620 and 640 TFLOPS for FP16 forward passes with a head dimension of 128 and sequence length of 8192).

\paragraph{Attention variants} In the cases of multi-query attention~\citep{shazeer2019fast} and grouped query attention~\citep{ainslie2023gqa}, our implementation aligns with the method used in \textsc{FlashAttention-2}, modifying tensor indexing strategies so that copies of $\mathbf{K}$ and $\mathbf{V}$ are not replicated in HBM.

\subsection{Intra-warpgroup overlapping GEMMs and softmax}

Within a single warpgroup, it is possible to concurrently execute certain instructions from the softmax computation alongside instructions from the GEMMs. Here, we outline a method to achieve this overlap.

In the attention algorithm, operations within the inner loop (main loop) have sequential dependencies that impede parallelization within a single iteration. For example, (local) softmax (lines \ref{code-ws:softmax_start} to \ref{code-ws:softmax_end}) relies on the output $\mathbf{S}_i^{(j)}$ of the first GEMM, while the second GEMM takes its result $\widetilde{\mathbf{P}}_i^{(j)}$ as an operand. Indeed, the wait statements in lines \ref{alg:ws_only_gemm1} and \ref{alg:ws_only_gemm2} of \cref{alg:flash3_wgmma_ws_only} serialize the execution of softmax and GEMMs. However, we can break these dependencies by pipelining across iterations through additional buffers in registers. Pursuing this idea, we propose the following two-stage\footnote{Note that the number of stages of the overlapping scheme is bounded by, but need not equal, the number $s$ of stages in the circular SMEM buffer.} GEMM-softmax pipelining algorithm:

\begin{algorithm}[H]
    \caption{\small\label{alg:flash3_wgmma}\textsc{FlashAttention-3} consumer warpgroup forward pass}
    \begin{algorithmic}[1]
      \REQUIRE  Matrices $\mathbf{Q}_i \in \mathbb{R}^{B_r \times d}$ and $\mathbf{K}, \mathbf{V} \in \mathbb{R}^{N \times d}$ stored in HBM, key block size $B_c$, with $T_c = \lceil \frac{N}{B_c} \rceil$.
      \STATE Reassign a set quantity of registers depending on the number of consumer warps utilized.
      \STATE Initialize $\mathbf{O}_i = (0) \in \mathbb{R}^{B_r \times d}$, as well as $\ell_i, m_i = (0), (-\infty) \in \mathbb{R}^{B_r}$ within on-chip memory.
      \STATE Synchronize until $\mathbf{Q}_i$ and $\mathbf{K}_0$ are available in shared memory.
      \STATE Using WGMMA, calculate $\mathbf{S}_{\mathrm{cur}} = \mathbf{Q}_i \mathbf{K}_0^T$. Finalize and await completion.
      \STATE Deallocate the initial stage in the buffer for $\mathbf{K}$.
      \STATE Determine $m_{i}$, $\tilde{\mathbf{P}}_{\mathrm{cur}}$, and $\ell_{i}$ from $\mathbf{S}_{\mathrm{cur}}$, then adjust $\mathbf{O}_i$ accordingly.
      \FOR{$1 \le j < T_c - 1$}
        \STATE Wait until $\mathbf{K}_j$ has been loaded into shared memory.
        \label{code:mainloop_start}
        \STATE Use WGMMA to compute $\mathbf{S}_{\mathrm{next}} = \mathbf{Q}_i \mathbf{K}_{j}^T$. Initiate commit process without waiting. \label{code:first_wgmma}
        \STATE Pause until $\mathbf{V}_{j-1}$ is ready in shared memory.
        \STATE Evaluate $\mathbf{O}_{i} = \mathbf{O}_{i} + \tilde{\mathbf{P}}_{\mathrm{cur}} \mathbf{V}_{j-1}$ with WGMMA, committing and proceeding immediately. \label{code:second_wgmma}
        \STATE Wait for the completion of WGMMA operation computing $\mathbf{Q}_i \mathbf{K}_{j}^T$.
        \STATE Update $m_{i}$, $\tilde{\mathbf{P}}_{\mathrm{next}}$, and $\ell_{i}$ from $\mathbf{S}_{\mathrm{next}}$. \label{code:softmax}
        \STATE Wait for $\tilde{\mathbf{P}}_{\mathrm{cur}} \mathbf{V}_{j-1}$ computation to finish and then rescale $\mathbf{O}_i$.
        \STATE Free the $(j\,\%\,s)$th and $(j-1\,\%\,s)$th stages of buffers allocated for $\mathbf{K}$ and $\mathbf{V}$ respectively.
        \STATE Assign $\mathbf{S}_{\mathrm{next}}$ to $\mathbf{S}_{\mathrm{cur}}$.
        \label{code:mainloop_end}
      \ENDFOR \STATE Await the loading of $\mathbf{V}_{T_c - 1}$ into shared memory.
      \STATE Perform the computation
      $\mathbf{O}_{i} = \mathbf{O}_{i} + \tilde{\mathbf{P}}_{\mathrm{last}} \mathbf{V}_{T_c - 1}$ with WGMMA. Commit and synchronize.

\STATE Epilogue: Adjust $\mathbf{O}_{i}$ according to $m_{i}$. Determine $L_{i}$ using both $m_{i}$ and $\ell_{i}$.
Store $\mathbf{O}_{i}$ and $L_{i}$ in HBM, assigning them to the $i$-th segment of $\mathbf{O}$ and $L$, respectively.
\end{algorithmic}
\end{algorithm}

\cref{alg:flash3_wgmma} substitutes for the consumer segment found in \cref{alg:flash3_wgmma_ws_only}, thereby completing the \textsc{FlashAttention-3} algorithm in FP16 mode. Conceptually, WGMMA is employed to represent asynchronous GEMM throughout the process. In the primary loop (spanning lines \ref{code:mainloop_start} to \ref{code:mainloop_end}), the second invocation of WGMMA within iteration $j$ (refer to line \ref{code:second_wgmma}) is executed concurrently with the softmax computations from the subsequent iteration $j+1$ (see line \ref{code:softmax}).

Although the pipelined organization discussed above promises improved theoretical efficiency, its practical application comes with several important considerations:

\paragraph{Compiler reordering} Although the pseudocode assumes an idealized sequence of execution, in practice the compiler (NVCC) frequently rearranges instructions for optimization purposes. Such reordering may interfere with the deliberate interleaving of WGMMA and non-WGMMA operations, risking either altered behavior or reduced performance benefits. Examination of the generated SASS code, however, confirms that the compiler does in fact produce overlapping instructions as intended (see Section~\ref{sec:2-stage-sass}).

\paragraph{Register pressure} Achieving maximum efficiency necessitates minimizing register spilling. Nevertheless, the 2-stage pipeline model requires additional register allocation to preserve intermediate computation states and support transitions between pipeline stages. Notably, a supplementary $\mathbf{S}_{\mathrm{next}}$ buffer must be maintained in registers, resulting in an additional per-threadblock register footprint of $B_r \times B_c \times \text{sizeof}(\text{float})$. Such increased demand on registers may conflict with other strategies, such as adopting larger block sizes, which likewise tend to be register-intensive. Accordingly, a careful balancing of trade-offs should be informed by empirical profiling.

\paragraph{3-stage pipelining} Building upon the previously discussed 2-stage design, we introduce a 3-stage pipeline that overlaps a second WGMMA stage with the softmax computation. Although this approach can further enhance Tensor Core utilization, it amplifies register requirements due to the presence of an additional stage, complicating the decision between maximizing tile size and deepening the pipeline. Details regarding this 3-stage method and its empirical assessment are provided in ~\cref{sec:3-stage}.

\subsection{Low-precision with FP8}
\label{sec:algofp8}

\textbf{Efficiency: layout transformations.} Executing the forward pass of \textsc{FlashAttention-3} with FP8 precision introduces layout compatibility complications that are not present when working with FP16.

To begin, it is important to recognize that the input tensors $\mathbf{Q}$, $\mathbf{K}$, and $\mathbf{V}$ are generally stored with head dimension as the contiguous axis. However, in order to comply with the k-major requirement of FP8 WGMMA for the second GEMM, $\mathbf{V}$—specifically, its tiles that are transferred into SMEM—must be contiguous along the sequence length dimension. As the TMA load process does not permit altering the contiguous dimension, two possible approaches emerge: (1) transposing $\mathbf{V}$ within GMEM during preprocessing, or (2) performing an in-kernel transpose of each tile after it has been loaded to SMEM. For the first method, the transpose operation can be integrated either by (1a) combining it with the epilogue of a preceding operation, such as rotary embedding, or by (1b) invoking a dedicated preprocessing transpose kernel\footnote{A high-performance transpose kernel can approach the device’s peak bandwidth~\citep{colfax_cutlass_transpose_2024}.} to switch the contiguous axes between sequence length and head dimensions. Nonetheless, incorporating the epilogue fusion as in (1a) poses substantial challenges for typical library implementations, while the standalone kernel as in (1b) leads to excessive resource consumption, especially during memory-bound scenarios such as inference.

For FP8 \textsc{FlashAttention-3}, we select option (2). Within the kernel, the transpose operation leverages the LDSM (\verb|ldmatrix|) and STSM (\verb|stmatrix|) instructions, allowing a warp to cooperatively read data from SMEM into RMEM and to write from RMEM back to SMEM in 128-byte units.\footnote{According to the PTX specification, LDSM/STSM are intended to copy $8 \times 8$ matrices composed of 16-bit elements \cite[\S 9.7.13.4.15-16]{ptx}. By packing pairs of 8-bit values together, it is possible to use LDSM/STSM with FP8 data. However, their transpose modes are unable to separate packed 8-bit pairs, which requires performing certain register shuffling between the LDSM and STSM to achieve a true tile-level transpose; we omit those implementation specifics here.} These instructions not only optimize register utilization—enabling them to run in the producer warpgroup—but also permit the layout to be transposed during memory transfers. Additionally, after the initial iteration, we may orchestrate the transposition of the subsequent $\mathbf{V}$ tile to occur concurrently with the execution of the two WGMMAs responsible for the previous $\mathbf{V}$ tile and the current $\mathbf{K}$ tile.

Secondly, it is evident that, in contrast to FP16, the memory organization of an FP32 accumulator within FP8 WGMMA deviates from the arrangement used for operand A when stored in registers. These structural differences are illustrated in \cref{fig:rmem_accum} and \cref{fig:rmem_operand}, where register entries are displayed per thread following their respective orderings. Through the use of byte permute instructions, we are able to convert the accumulator from the initial WGMMA into a layout that aligns with what is required by the subsequent WGMMA and also matches the $\mathbf{V}$ tile structure generated via the in-kernel transpose. More concretely, as seen in \cref{fig:rmem_accum}, the register contents are reordered as follows:
$$\{ \verb|d0 d1 d4 d5 d2 d3 d6 d7| \},$$
with this pattern repeated every 8 bytes within the registers. This process effectively permutes the columns of the $\mathbf{P}$ tile; for instance, columns $0189$ are reassigned as the first set of columns. To ensure the second WGMMA yields the correct output tile, the in-kernel transpose can be implemented such that the rows in the $\mathbf{V}$ tile are permuted in a matching manner.\footnote{By leveraging the additional flexibility provided by the in-kernel transpose, it is unnecessary to use shuffle instructions for transferring register data across threads, as previously discussed in~\citep{colfax_fp8_flashattention_2024}.}

\textbf{Accuracy: block quantization and incoherent processing.} The FP8 (e4m3) representation allocates only 3 bits for the mantissa and 4 bits for the exponent, inherently producing greater numerical error relative to FP16/BF16. Additionally, large-scale models often contain outlier entries~\citep{dettmers2208llm, sun2024massive}—values with magnitudes that far exceed the majority—posing significant challenges for quantization. A common strategy is per-tensor scaling~\citep{micikevicius2022fp8}, where each tensor is assigned a single scalar scaling factor (such as one for $\mathbf{Q}$, one for $\mathbf{K}$, and one for $\mathbf{V}$). To address and reduce the precision loss encountered when executing attention in FP8, we introduce two complementary methods:

\begin{enumerate}

\item \textbf{Block quantization}: In this method, a single scalar per block is maintained. For each of $\mathbf{Q}$, $\mathbf{K}$, and $\mathbf{V}$, the tensor is partitioned into blocks of size $B_r \times d$ or $B_c \times d$ and each block is quantized independently. This quantization step can be seamlessly integrated with preceding operations, such as rotary embedding, without incurring any additional computational latency, as rotary embedding is limited by memory bandwidth. Given that the \textsc{FlashAttention-3} algorithm natively processes data in blocks, it is possible to adjust the scaling for each block in $\mathbf{S}$ to reflect the effects of block quantization, without any computational overhead.

\item \textbf{Incoherent processing}: To mitigate the effect of outliers, both $\mathbf{Q}$ and $\mathbf{K}$ are premultiplied by a random orthogonal matrix $\mathbf{M}$ prior to FP8 quantization. Since $\mathbf{M}$ satisfies $\mathbf{M} \mathbf{M}^\top = I$, the resulting product $(\mathbf{Q} \mathbf{M}) (\mathbf{K} \mathbf{M})^\top$ is equivalent to $\mathbf{Q} \mathbf{K}^\top$, ensuring that the attention computation remains unchanged. This process disperses outliers because each component in $\mathbf{Q} \mathbf{M}$ or $\mathbf{K} \mathbf{M}$ is a randomly weighted sum of the original entries in $\mathbf{Q}$ or $\mathbf{K}$, thereby decreasing quantization error. Following the approach in \citet{chee2024quip} and \citet{tseng2024quip}, we select $\mathbf{M}$ as the product of random diagonal matrices with entries $\pm 1$ and a Hadamard matrix, which allows multiplication to be performed in $O(d \log d)$ time, rather than $O(d^2)$, and this can also be merged with the rotary embedding without adding extra computation.

We demonstrate in \cref{sec:numerical_error} that these approaches can reduce numerical error by up to a factor of 2.6.

\section{Empirical Validation}
\label{sec:exp}

To implement \textsc{FlashAttention-3} and assess its performance and precision, we leverage CUTLASS~\citep{Thakkar_CUTLASS_2023} primitives, including the WGMMA and TMA abstractions.

\begin{itemize}

\item \textbf{Benchmarking attention.} The execution time of \textsc{FlashAttention-3} was evaluated over a range of sequence lengths and contrasted with a baseline PyTorch implementation, \textsc{FlashAttention-2}, \textsc{FlashAttention-2} in Triton (utilizing H100-targeted instructions), as well as a commercial release of \textsc{FlashAttention-2} specifically optimized for H100 GPUs within cuDNN. Results indicate that \textsc{FlashAttention-3} achieves up to a 2.0$\times$ speedup relative to \textsc{FlashAttention-2} and a 1.5$\times$ acceleration compared to the Triton implementation. The throughput of \textsc{FlashAttention-3} reaches up to 740 TFLOPs/s, which is approximately 75\% of the H100 GPU's peak theoretical TFLOPs/s. 

\item \textbf{Ablation study.} Through experimentation, it is demonstrated that advances such as warp-specialization and the use of GEMM-softmax pipelining significantly drive the performance improvements observed in \textsc{FlashAttention-3}.

\item \textbf{Accuracy of FP8 attention.} Our assessment confirms that integrating block quantization alongside incoherent processing effectively decreases the numerical error associated with FP8 \textsc{FlashAttention-3} by a factor of 2.6$\times$.
\end{itemize}

\subsection{Benchmarking Attention}
\label{subsec:benchmark_attn}

We evaluate the execution time of various attention mechanisms on an H100 80GB SXM5 GPU under multiple configurations (with and without causal masking, head sizes of 64 and 128), using FP16 precision. The outcomes are presented in~\cref{fig:benchmark_attn_fwd} and~\cref{fig:benchmark_attn_bwd}. These results reveal that \textsc{FlashAttention-3} achieves a speedup of approximately 1.5-2.0$\times$ over \textsc{FlashAttention-2} in the forward direction, and is about 1.5-1.75$\times$ faster in the backward direction. When compared to a conventional implementation of attention, \textsc{FlashAttention-3} demonstrates accelerations ranging from 3$\times$ to 16$\times$. Notably, for sequences of medium and longer lengths (1k tokens or more), \textsc{FlashAttention-3} is able to outperform even the highly optimized proprietary cuDNN library for the H100 GPU.

\paragraph{Benchmark settings:} We vary the sequence length as 512, 1k, ..., 16k, and set batch size so that the total number of tokens is 16k. We set the hidden dimension to 2048, and head dimension to be either 64, 128, or 256 (i.e., 32 heads, 16 heads, or 8 heads). To calculate the FLOPs of the forward pass, we use:

\begin{equation*}
  4 \cdot \text{seqlen}^2 \cdot \text{head dimension} \cdot \text{number of heads}.
\end{equation*}

By applying causal masking, we reduce this value by half, since roughly 50% of the entries are actually computed. The FLOPs required for the backward pass are then estimated by multiplying the forward pass FLOPs by 2.5, reflecting the difference in matrix multiplications: the forward pass involves 2 matmuls, whereas the backward pass requires 5 matmuls owing to recomputation.

Additionally, we assess the forward pass runtime using FP8 in comparable conditions. The outcomes for headdim 256 are presented in ~\cref{fig:benchmark_attn_fp8}, with comprehensive results detailed in \cref{sec:benchmark_attn_fp8_full}.

\subsection{Ablation Study: 2-Stage Pipelining Experiments}
\label{sec:2-stage-pipelining-experiments}

We conduct an ablation study on both the 2-stage WGMMA-softmax pipeline and warp-specialization for non-causal FP16 \textsc{FlashAttention-3}, keeping the parameters fixed at $\{\text{batch}, \text{seqlen}, \text{nheads}, \text{hdim}\} = \{ 4, 8448, 16, 128\}$. The outcomes presented in~\cref{table:ablation_pipelining} demonstrate that the combination of asynchrony through warp-specialization and the concurrent execution of GEMM and softmax substantially accelerates performance, achieving a throughput increase from 570 to 661 TFLOPs.

\subsection{Numerical Error Validation}
\label{sec:numerical_error}

Given increased scrutiny regarding the numerical error~\citep{golden2024flash} associated with \textsc{FlashAttention}, we evaluate \textsc{FlashAttention-2}, \textsc{FlashAttention-3}, and a conventional attention mechanism by benchmarking them against a reference implementation performed in FP64 precision. In order to emulate atypical features and activations present in LLMs~\citep{dettmers2208llm, sun2024massive}, the elements of $\mathbf{Q}, \mathbf{K}, \mathbf{V}$ are sampled according to the following distribution:

\begin{equation*}
  \mathcal{N}(0, 1) + \mathcal{N}(0, 100) \cdot \mathrm{Bernoulli}(0.001).
\end{equation*}

Specifically, the elements are sampled from a normal distribution with a mean of zero and a standard deviation of 1; however, for 0.1\% of the elements, an additional independent normal component with standard deviation 10 is included. The root mean squared error (RMSE) is reported in~\cref{table:numerical_error}. When using FP16, \textsc{FlashAttention-2} and \textsc{FlashAttention-3} yield RMSE values that are 1.7$\times$ lower than those from the standard implementation, since softmax computations are maintained in FP32. For the baseline FP8 attention, per-tensor scaling is adopted, matrix multiplication accumulations occur in FP32, and intermediate softmax data is stored in FP16. Owing to the utilization of block quantization and incoherent processing, \textsc{FlashAttention-3} operating in FP8 attains an accuracy 2.6$\times$ greater than this baseline approach.

\section{Dicussion, Limitations, Conclusion}
\label{sec:discussion}

Through \textsc{FlashAttention-3}, we illustrate that leveraging innovative programming paradigms and modern hardware capabilities, including asynchrony and reduced precision, can substantially enhance both the efficiency and precision of attention mechanisms. Compared to \textsc{FlashAttention-2}, attention is accelerated by a factor of 1.5–2.0$\times$, and FP8 numerical errors are mitigated by 2.6$\times$ relative to traditional per-tensor quantization approaches. Our current work faces certain challenges that remain to be solved: refining the approach for LLM inference scenarios, implementing a persistent kernel within the FP8 kernel,\footnote{For the benchmarking setup, FP16 \textsc{FlashAttention-3} includes both a persistent kernel and an effective load-balancing method, whereas FP8 \textsc{FlashAttention-3} does not, contributing to its weaker performance for short sequence lengths and causal masks in comparison with FP8 cuDNN kernels.} and investigating the influence of low-precision attention during large-scale training regimes. Although the present study centers on Hopper GPUs, we anticipate these strategies will be extendable to a variety of hardware accelerators. Ultimately, we foresee that the improved speed and fidelity of this attention primitive will enable novel solutions for tasks requiring extended context.

\end{document}